\label{chap:introduction}

A recently discovered class of hardware side-channel vulnerabilities
allows an attacker to recover remnants of prior activity from transistor
devices in FPGA routing logic long after their logical erasure
\cite{drewes2024pentimento}. The remnants that these attacks exploit
have been termed \emph{pentimenti}, after the art history term for
traces of an earlier composition that remain beneath the finished
surface of a painting. As conventional pentimenti are revealed through
X-ray radiography or infrared reflectography \cite{gooch2007visualizing,
	daffara2011multispectral}, FPGA pentimenti are revealed through
fine-grained timing measurement. These measurements reveal subtle
variations in propagation delay that reveal physical evidence of the
device's prior state, and prior users' activity.

To date, FPGA pentimenti have been demonstrated only in the
reconfigurable routing logic of Xilinx field-programmable gate arrays
(FPGAs). The underlying mechanism is bias temperature instability (BTI),
a recoverable wearout phenomenon that causes small, persistent drifts in
transistor switching time under a prolonged stress bias. When a single
bit of information is held in an FPGA routing net for an extended
period, the accumulated BTI stress across all transistors propagating
that bit produce subtle delay shifts large enough to be resolved by
time-to-digital converter (TDC) sensors implemented in soft logic. The
accuracy of these TDCs depends on phase-shift granularity provided by
the FPGA's phase-locked loops (PLLs), which varies considerably between
vendors. Prior demonstrations of FPGA pentimenti recovery have relied on
the high-resolution PLL phase-shift capabilities available only on
Xilinx products. As a first contribution, this dissertation closes this
gap. We reformulate fine-grained sensor tuning as a constrained PLL
configuration search, incorporating both phase shifting and frequency
synthesis to compose a fine-grained phase sweep on Intel Cyclone V FPGAs
without relying on high-resolution vendor-specific phase-shift
capabilities.

With pentimenti measurement now broadly accessible across FPGA vendors,
we ask the following question: does the FPGA pentimenti attack class
extend beyond FPGAs? BTI affects every CMOS transistor to varying
degrees \cite{hill2021cmos}, and analogous defect dynamics do occur in
many other semiconductor devices such as self-recovery in 3-D NAND cells
\cite{luo2018heatwatch} and relaxation in phase-change memory
\cite{le2018collective}. However, no techniques for constructing a
non-invasive leakage channel like that of FPGA pentimenti attacks have
been proposed.

Resistive random-access memory (ReRAM) is a compelling candidate
technology for the study of pentimenti-based attacks. In ReRAM, each bit
is physically stored as a conductive filament state whose growth and
rupture is dictated by stochastic defect dynamics~\cite{waser2009redox,
	ambrogio2013understanding, Ielmini_2016}. The duration of each
programming operation is dictated by these processes
\cite{jo2009programmable, witzleben2017investigation}. This allows for
leakage channels to be established which leak defect-dynamics-driven
behavior through measurements of write latency. Unlike FPGAs, however,
commercial ReRAM modules expose only a memory interface and defect
dynamics cannot be measured directly as in FPGA pentimenti. Instead, the
observable write latency reflects word-level switching activity that is
dominated by data-dependent structural biases. These biases arise from
write-controller polarity serialization across cells within a word and
additional cell transitions introduced by undocumented embedded
error-correcting codes (ECCs). Isolating the underlying stochastic
component of programming time requires models of these structural timing
biases.

This dissertation develops the reverse engineering and measurement
infrastructure necessary to study pentimenti-class side-channels in
commercial ReRAM. We characterize the dominant structural sources of
write latency in commercial ReRAM crossbars, introduce \emph{Bit-exact
	Recovery of ECC from Write Timing in ReRAM Crossbars} (BREW-RC) as the
first non-invasive technique for recovering an embedded ECC parity
submatrix from write timing alone, and use this structural understanding
to isolate and characterize the stochastic component of write latency.
Together, these contributions establish the required measurement
capabilities for ReRAM pentimenti analysis. The remainder of this
chapter states the research objectives and contributions of this
dissertation and describes the organization of the remaining chapters.